% !TEX root = IE3_expI_prelab.tex
\documentclass[lualatex,a4paper,10pt]{jlreq}
% Shared LuaLaTeX preamble.
\usepackage{luatexja-fontspec}
\setmainfont{TeX Gyre Termes}
\setsansfont{TeX Gyre Heros}
\setmonofont{Latin Modern Mono}
\setmainjfont[BoldFont=HaranoAjiGothic-Medium]{HaranoAjiMincho-Regular}
\setsansjfont[BoldFont=HaranoAjiGothic-Bold]{HaranoAjiGothic-Medium}
\usepackage[top=22mm,bottom=22mm,left=23mm,right=23mm]{geometry}
\usepackage{amsmath,amssymb,booktabs,array,longtable,tabularx}
\usepackage{graphicx,xcolor,tikz,circuitikz}
\usetikzlibrary{arrows.meta,positioning,calc}
\usepackage{enumitem,fancyhdr,listings,needspace}
\usepackage[unicode,colorlinks=true,linkcolor=labblue,urlcolor=labteal]{hyperref}
\usepackage{bookmark}
\definecolor{labblue}{RGB}{30,70,112}
\definecolor{labteal}{RGB}{0,100,105}
\definecolor{labgray}{gray}{0.95}
\ModifyHeading{section}{font={\large\sffamily\bfseries\color{labblue}}}
\ModifyHeading{subsection}{font={\normalsize\sffamily\bfseries\color{labteal}}}
\setlength{\parindent}{1\zw}
\setlength{\parskip}{3pt}
\setlength{\emergencystretch}{2em}
\renewcommand{\arraystretch}{1.35}
\setlist{itemsep=3pt,topsep=4pt,leftmargin=2\zw}
\newcolumntype{Y}{>{\raggedright\arraybackslash}X}
\newcolumntype{C}[1]{>{\centering\arraybackslash}p{#1}}
\pagestyle{fancy}
\fancyhf{}
\fancyhead[L]{\small\color{labblue} IE3 後期・電子工学実験}
\fancyhead[R]{\small テーマ\LabCode}
\fancyfoot[C]{\thepage}
\setlength{\headheight}{16pt}
\DeclareOldFontCommand{\rm}{\normalfont\rmfamily}{\mathrm}
\lstset{basicstyle=\small\ttfamily,columns=fullflexible,keepspaces=true,
 breaklines=true,frame=single,backgroundcolor=\color{labgray},
 numbers=left,numberstyle=\tiny,showstringspaces=false,aboveskip=8pt,belowskip=8pt}
\newcommand{\blank}{\rule{22mm}{0.3pt}}
\newcommand{\answerline}{\par\noindent\rule{\linewidth}{0.25pt}\par}
\newcommand{\writing}[1]{\par\Needspace{\dimexpr#1+5mm\relax}\noindent%
 \fcolorbox{labblue!35}{white}{\parbox[c][#1][t]{\dimexpr\linewidth-2\fboxsep-2\fboxrule\relax}{\vspace{2mm}\noindent\strut\hfill\par}}\par}
\newcommand{\hint}[1]{\par\smallskip\noindent{\sffamily\bfseries\color{labteal}記述の視点：}#1\par}
\newcommand{\note}[1]{\par\smallskip\noindent\fcolorbox{labblue!45}{labblue!5}{%
 \parbox{\dimexpr\linewidth-2\fboxsep-2\fboxrule\relax}{#1}}\par\smallskip}
\newcommand{\eqerror}{\varepsilon=\frac{x_{\mathrm{meas}}-x_{\mathrm{th}}}{x_{\mathrm{th}}}\times100\ \%}
\newcommand{\LabSource}{徳山工業高等専門学校 情報電子工学科，
 『2026年度 電子工学実験（3年後期）テキスト』，テーマ\LabCode，
 \expandafter\path\expandafter{\SourcePdf}。}
\newcommand{\reporttitle}{
 \noindent{\small\sffamily\color{labteal}実験報告書・記入ガイド}\par
 \noindent{\LARGE\sffamily\bfseries \LabCode\quad\LabTitle}\par\smallskip
 \noindent 氏名：\blank\quad 班：\blank\quad 実験日：\blank\par
 \note{目的・原理・手順を確認し、実測値と自分の考察を記入する。
 考察のガイドは、検討する事象・使う測定結果・記述の方針を示す。}
}
\newcommand{\prelabtitle}{
 \noindent{\small\sffamily\color{labteal}事前学習・前提知識プリント}\par
 \noindent{\LARGE\sffamily\bfseries \LabCode\quad\LabTitle}\par\smallskip
 \noindent 氏名：\blank\quad 班：\blank\quad 予習日：\blank\par
 \note{用語、回路の動き、計算、測定表への記入の順に読む。
 例題の値は説明用であり、実験結果ではない。}
}
\newcommand{\tocrow}[3]{\hyperref[#1]{#2}& #3 &\pageref*{#1}\\[2mm]}
\newcommand{\documentmap}[1]{
 \section*{目次と概要}
 \noindent 青色の項目をクリックすると該当箇所へ移動する。\par
 {\small\begin{longtable}{@{}p{0.29\linewidth}p{0.61\linewidth}r@{}}
 \toprule {\color{labblue}\bfseries 項目} & {\color{labblue}\bfseries 読むと分かること} & 頁\\\midrule
 \endhead #1\bottomrule
 \end{longtable}}
 \clearpage
}
\newenvironment{terms}{\par\smallskip\begin{longtable}{@{}p{0.23\linewidth}p{0.73\linewidth}@{}}
 \toprule {\color{labteal}\bfseries 用語}&{\color{labteal}\bfseries 意味・回路での役割}\\\midrule\endhead}
 {\bottomrule\end{longtable}}
\newcommand{\term}[2]{\textbf{#1}&#2\\[2mm]}
\newcommand{\guidepoint}[4]{
 \Needspace{32mm}\subsection{#1}
 \noindent{\bfseries\color{labteal}考える事象：}#2\par
 \noindent{\bfseries\color{labteal}参照する結果：}#3\par
 \noindent{\bfseries\color{labteal}記述の方針：}#4\par
}
\newcommand{\reportclosing}[1]{
 \clearpage\section{結論のまとめ方}\label{sec:closing}
 \hint{#1}
 \ConclusionGuide
 \begin{enumerate}
 \item 目的に対応する最も重要な実測結果を、測定条件と数値を添えて述べる。
 \item 原理と一致した点・一致しなかった点を示す。原因は確認済みの事実と推測を分ける。
 \item 実施範囲で言えることと、未確認の条件を一文ずつまとめる。
 \end{enumerate}
 \note{書き出しの型：「○○の条件で△△を測定したところ、□□となった。
 これは◇◇と整合する。ただし××は確認しておらず、一般化には追加測定が必要である。」
 空欄は自分の実測値と根拠で埋める。}
 \noindent 結論（実験後に記入）：\writing{30mm}
 \noindent 改善案・次に確認したい条件：\writing{16mm}
 \section*{参考文献}\label{sec:references}
 \begin{enumerate}
 \item \LabSource
 \item 使用したデータシート・書籍・ウェブ資料（著者、題名、該当箇所、URL、参照日）を追加する。
 \end{enumerate}
}
\newcommand{\prelabclosing}{
 \section*{準備・出典}\label{sec:prelabsource}
 \begin{itemize}[nosep]
 \item 資料、筆記具、記録用紙、電卓と、実験条件・機器設定の記録欄。
 \item 確認問題の解答と、分からない用語・回路点への具体的な質問。
 \end{itemize}
 {\small\noindent\LabSource\par}
}
\newcommand{\simpleflow}[1]{
 \begin{center}
 \begin{tikzpicture}[node distance=5mm and 5mm,
 every node/.style={font=\small},box/.style={draw=labblue,fill=labblue!4,rounded corners=1mm,
 align=center,minimum height=10mm,text width=30mm}]
 #1
 \end{tikzpicture}
 \end{center}
}

\newcommand{\LabCode}{I}
\newcommand{\LabTitle}{A/D変換とD/A変換}
\newcommand{\SourcePdf}{IE3_expI_A4.pdf}
\begin{document}
\prelabtitle
\documentmap{
\tocrow{sec:auto1}{1 用語を一つずつ理解する}{言葉の意味、単位、回路や測定での役割を確認する。}
\tocrow{sec:auto2}{2 三つの変換回路をたどる}{部品と信号の経路を追い、状態が変化する理由を理解する。}
\tocrow{sec:auto3}{3 量子化と符号}{ビット数、LSB、最大出力と誤差を読む。}
\tocrow{sec:auto4}{4 R--2Rと変換方式}{この項目の仕組みと実験で確認すべき点を整理する。}
\tocrow{sec:auto5}{5 サンプリングの前提}{標本間隔と信号再構成の条件を理解する。}
\tocrow{sec:auto6}{6 折返しの例}{入力と標本化周波数から見かけの周波数を求める。}
\tocrow{sec:auto7}{7 サンプルホールド}{取得と保持の期間を波形で区別する。}
\tocrow{sec:auto8}{8 測定からレポートへ：符号と波形を読む}{実測値から計算・作図・記述へ進む順序を例題で練習する。}
\tocrow{sec:auto9}{9 確認問題}{自分で計算・説明して、実験に必要な理解を確かめる。}
\tocrow{sec:auto10}{解答の目安}{数値と理由を照合し、単位や論理の取り違えを見直す。}
\tocrow{sec:auto11}{補足資料}{メーカー公式資料で方式・端子・動作条件を確認する。}
\tocrow{sec:prelabsource}{準備・出典}{持参物と資料を確認し、具体的な質問を準備する。}
}

\section{用語を一つずつ理解する}\label{sec:auto1}
\begin{terms}
\term{アナログ・デジタル}{連続する電圧を扱うのがアナログ、離散した符号を扱うのがデジタル。デジタル化しても測定誤差が消えるわけではない。}
\term{ビット・MSB・LSB}{nビットには$2^n$通りの符号がある。MSBは最も大きい重みのビット、LSBは最小のビットであり、一段分の電圧の意味でも使う。}
\term{量子化・分解能}{連続電圧を段階へ対応させる操作と、その最小刻み。段階が増すほど細かく表せるが、絶対精度は基準・素子誤差にも依存する。}
\term{基準電圧・フルスケール}{符号と電圧のものさし。資料の8ビットDACは2.5 V基準で変換後、増幅回路で2倍にするので、理想の出力スケールは5 Vになる。}
\term{比較器・エンコーダ}{大小を判定する回路と、判定結果を2進符号へ変える回路。優先エンコーダは複数の有効入力から優先順位を選ぶ。}
\term{標本化・折返し}{時間を離散的に取り出す操作と、不十分な標本化で別の周波数に見える現象。電圧の量子化と時間の標本化は別の制約である。}
\term{サンプルホールド}{取得区間で入力へ追従し、保持区間では取得した電圧を保つ回路。取得時間や漏れによる変化もある。}
\end{terms}
\section{三つの変換回路をたどる}\label{sec:auto2}
\subsection{3ビット並列比較形ADC}
基準を分けた複数のしきい値と入力を、比較器で同時に比べる。
入力が高くなると、越えたしきい値の数が増える。その並んだ判定を優先エンコーダで3ビットへ変換する。
74LS148は入力と出力がLOWで有効となるICなので、端子の生の0/1と「入力の大きさを表す符号」を無条件に同一視しない。資料の反転・配線とLED表示まで追う。
\subsection{R--2R DAC：各ビットの寄与を足す}
各ビットは0 Vか基準電圧を選ぶスイッチとして働く。
Rと2Rが繰り返される構造により、寄与が上位から1/2、1/4、1/8となる。
一つのビットだけを1にして他を0にすると、そのビットの寄与を調べられる。これを重ねの理で足す。
出力のバッファは後段の負荷が抵抗ラダーへ影響するのを抑える役割を持つ。
\clearpage
\subsection{8ビットICと出力増幅}
ADC0804は逐次比較形で、上位ビットから仮のDAC出力を入力と比較して符号を決める。
変換には時間が必要であり、標本化の開始と全ビット確定は同時とは限らない。
資料のDACはTLC7524で、ラダー変換の後に非反転増幅回路がある。
ADCの入力電圧、DAC単体の電圧、最終出力電圧を区別する。
\subsection{サンプルホールドの二状態}
取得中はスイッチを通して容量の電圧を入力へ近づけ、保持中はスイッチを離して値を保つ。
バッファは保持容量から負荷へ電荷が流れにくくする。
取得がHIGHかLOWかは資料のスイッチ制御を確認する。100 Hz入力と50 Hz取得では同位相を繰り返し取る場合があるため、保持値が変わらなくても直ちに故障と判断しない。

\section{量子化と符号}\label{sec:auto3}
A/Dは連続電圧を離散符号へ、D/Aは符号をアナログ値へ変換する。
5 Vスケールで3ビットなら$2^3=8$段階、LSBは0.625 V。
8ビットなら256段階、LSBは約19.53 mVである。
理想的な四捨五入量子化の誤差は最大半LSBだが、実際の符号境界・オフセットは方式とデータシートに従う。
\[
 V_o=5\frac D{2^n},\quad V_{o,\max}=5\left(1-\frac1{2^n}\right).
\]
最大値は3ビットで4.375 V、8ビットで4.98047 V。
\section{R--2Rと変換方式}\label{sec:auto4}
R--2Rラダーは2種類の抵抗で重み付き和を作る。
3ビットなら$V_o=5(b_2/2+b_1/4+b_0/8)$と表せる。
入力101は3.125 Vになる。
並列比較形は複数しきい値を同時比較する。逐次比較形はDACと比較器で上位ビットから判定する。
資料の3ビット回路は前者、ADC0804は後者である。
\section{サンプリングの前提}\label{sec:auto5}
標本間隔$T_s=1/f_s$。入力一周期の標本数は$f_s/f_{\rm in}$。
理想的に帯域制限された信号の再構成条件は$f_s>2f_{\max}$。
実回路ではアンチエイリアスフィルタ、変換時間、保持、再構成フィルタも関係する。
\clearpage
\section{折返しの例}\label{sec:auto6}
観測帯域$0$〜$f_s/2$に入るよう整数kを選んだ
\[
 f_{\rm alias}=|f_{\rm in}-kf_s|
\]
が折返し周波数になる。
50 Hz入力を60 Hzで標本化すると10 Hz、110 Hz入力を80 Hzで標本化すると30 Hzに見える。
階段波形を単に滑らかに見せることと、元信号を復元できることは異なる。
\section{サンプルホールド}\label{sec:auto7}
スイッチがONの取得期間でコンデンサを入力に追従させ、OFFの保持期間で値を保持する。
100 Hz入力と50 Hz・20\%デューティのパルスなら、パルス周期20 ms、指定レベルの期間4 msとなる。実際の取得極性は回路で確認する。
\section{測定からレポートへ：符号と波形を読む}\label{sec:auto8}
\subsection{2進・16進から電圧へ}
3ビット101は$4+0+1=5$で、5 Vスケールなら$5\times5/8=3.125\,\mathrm{V}$。
各寄与でも$2.5+0+0.625=3.125\,\mathrm{V}$となる。
8ビット80hは128、FFhは255。最大出力は$5\times255/256$であり5 Vちょうどではない。
ADCの同じ符号には電圧範囲が対応するので、入力は一点だけでなく切り替わる境界も記録する。
\subsection{ビットを一つずつ立てる理由}
01h、02h、04h等の出力がほぼ倍々になれば、各ビットの重みを確認できる。
零符号のずれがあればオフセット、値に比例して増すずれがあればスケールの差を候補とする。
この区別には一つの最大値だけでなく、複数符号の測定が必要である。
\subsection{標本数と折返し}
入力50 Hz、標本化1000 Hzなら一周期20標本。
50 Hzを60 Hzで標本化すると10 Hz、110 Hzを80 Hzでは30 Hzの見かけの変化を生じ得る。
標本間の段差と、一周期全体が別の周波数に見える折返しを分けて記す。
「一周期10点で見やすい」という目安と、帯域制限された信号の再構成条件$f_s>2f_{\max}$は意味が異なる。
\subsection{保持波形の記録}
入力・制御パルス・出力の同じ時刻に線を引き、追従する区間と水平に近い区間を対応させる。
出力がずっと平らな場合は、取得周期と入力位相、回路設定、スコープの結合を確認する。
\par\smallskip
74LS148の論理極性の補足：Texas Instruments公式資料，
\url{https://www.ti.com/product/SN74LS148}（参照日：2026-10-09）。

\clearpage
\section{確認問題}\label{sec:auto9}
\begin{enumerate}
\item 3ビットの110の理想DAC出力を求める。\answerline
\item 8ビットの80hとFFhの理想出力を求める。\answerline
\item 入力50 Hz、サンプリング500 Hzの一周期あたり標本数を求める。\answerline
\item 110 Hzを80 Hzで標本化すると、再構成条件を満たすか。折返し周波数も答える。\answerline
\end{enumerate}
\section*{解答の目安}\label{sec:auto10}
(1) 3.75 V。(2) 2.5 V、約4.9805 V。(3) 10。
(4) 満たさない、30 Hz。
\section*{補足資料}\label{sec:auto11}
Texas Instruments, \textit{ADC080x 8-Bit, $\mu$P-Compatible ADCs}：
\url{https://www.ti.com/lit/ds/symlink/adc0804-n.pdf}
（参照日：2026-10-09）。
\prelabclosing
\end{document}
